% ============================================================
%  Soma dos Ângulos de um Triângulo — Apostila Premium | PratiqueJá
%  Engine: XeLaTeX
% ============================================================
\documentclass[12pt]{article}
\usepackage{../../pratiqueja}
\usepackage{tikz}
\usetikzlibrary{calc}

% \hlm — destaque de resultado em modo-math (cor sem bold)
\newcommand{\hlm}[1]{{\color{darkblue}#1}}

\tikzset{
  tri/.style={draw=darkblue, line width=1.2pt, fill=accentblue!12!white},
  aalpha/.style={vered, line width=1.2pt},
  abeta/.style={badgepurple, line width=1.2pt},
  agamma/.style={badgegreen, line width=1.2pt},
  cl/.style={font=\footnotesize\bfseries},
}

\setsubject{Soma dos Ângulos de um Triângulo}

\begin{document}
\pagenumbering{arabic}
\setstretch{1.2}
\color{bodytext}

\heroheader{Ângulos do Triângulo}%
           {Soma dos Internos, Isósceles e Ângulo Externo}%
           {Matemática · Intermediário}

\setlength{\columnsep}{18pt}
\setlength{\columnseprule}{0.5pt}
\def\columnseprulecolor{\color{exborder}}

\begin{multicols}{2}

%─────────────────────────────────────────────────────────────
\TW{Def}{badgepurple}{Definição}{Soma dos ângulos internos}

\regra{A soma dos ângulos internos de \textbf{qualquer triângulo} é
sempre igual a $180°$.}

\formula{$\alpha + \beta + \gamma = 180°$}

\begin{center}\vspace{2pt}
\begin{tikzpicture}[scale=0.82]
\coordinate (A) at (0,0); \coordinate (B) at (5,0); \coordinate (C) at (1.5,2.5);
\draw[tri] (A)--(B)--(C)--cycle;
\draw[aalpha]  (A)++(0:0.55) arc (0:59:0.55);
\draw[agamma]  (B)++(145:0.55) arc (145:180:0.55);
\draw[abeta]   (C)++(239:0.55) arc (239:324:0.55);
\node[vered,cl]      at (0.66,0.30){$\alpha$};
\node[badgegreen,cl] at (4.22,0.22){$\gamma$};
\node[badgepurple,cl] at (1.5,1.78){$\beta$};
\end{tikzpicture}
\end{center}

\obs{Vale para escaleno, isósceles ou equilátero. Serve para achar
ângulos desconhecidos e fundamenta identidades trigonométricas.}

\SH{Exemplo — achar o 3º ângulo}
\exemp{%
  Ângulos de $65°$ e $30°$, achar $x$:\\
  $65° + 30° + x = 180°$\\
  $95° + x = 180° \Rightarrow x = 180° - 95° = \hlm{85°}$%
}

%─────────────────────────────────────────────────────────────
\TW{Iso}{badgegreen}{Triângulo Isósceles}{Dois lados e dois ângulos iguais}

\regra{No triângulo \textbf{isósceles}, os dois ângulos da \textbf{base}
são sempre iguais.}

\begin{center}\vspace{2pt}
\begin{tikzpicture}[scale=0.78]
\coordinate (A) at (0,0); \coordinate (B) at (3,0); \coordinate (C) at (1.5,3);
\draw[tri] (A)--(B)--(C)--cycle;
\draw[aalpha] (A)++(0:0.55) arc (0:63:0.55);
\draw[aalpha] (B)++(117:0.55) arc (117:180:0.55);
\draw[abeta]  (C)++(244:0.5) arc (244:296:0.5);
% marcas de lados congruentes
\draw[darkblue,line width=1pt] ($(0.75,1.5)+(153:0.12)$)--($(0.75,1.5)+(153:-0.12)$);
\draw[darkblue,line width=1pt] ($(2.25,1.5)+(27:0.12)$)--($(2.25,1.5)+(27:-0.12)$);
\node[vered,cl]       at (0.62,0.32){$\alpha$};
\node[vered,cl]       at (2.38,0.32){$\alpha$};
\node[badgepurple,cl] at (1.5,2.18){$\beta$};
\end{tikzpicture}
\end{center}

\obs{Possui \textbf{dois lados de mesma medida}; o terceiro é a
\textbf{base}. Os ângulos entre a base e os lados congruentes são iguais.}

\SH{Exemplo — ângulo do vértice $40°$}
\exemp{%
  Ângulos da base $x$ iguais, vértice $40°$:\\
  $x + x + 40° = 180° \Rightarrow 2x = 140°$\\
  $x = \dfrac{140°}{2} = \hlm{70°}$%
}

%─────────────────────────────────────────────────────────────
\TW{Ext}{darkblue}{Ângulo Externo}{Igual à soma dos internos não adjacentes}

\regra{O \textbf{ângulo externo} é igual à \textbf{soma dos dois ângulos
internos não adjacentes} a ele.}

\formula{$\text{externo} = \alpha + \beta$}

\begin{center}\vspace{2pt}
\begin{tikzpicture}[scale=0.78]
\coordinate (A) at (0,0); \coordinate (B) at (5,0); \coordinate (C) at (1.5,2.5);
\coordinate (X) at (6.5,0);
\draw[tri] (A)--(B)--(C)--cycle;
\draw[line width=1.2pt, darkblue] (B)--(X);
\draw[aalpha] (A)++(0:0.55) arc (0:59:0.55);
\draw[abeta]  (C)++(239:0.55) arc (239:324:0.55);
\draw[agamma] (B)++(145:0.5) arc (145:180:0.5);
\draw[vered, line width=1.2pt] (B)++(0:0.6) arc (0:145:0.6);
\node[vered,cl]       at (0.66,0.30){$\alpha$};
\node[badgepurple,cl] at (1.5,1.78){$\beta$};
\node[badgegreen,cl]  at (4.3,0.24){$\gamma$};
\node[vered,cl]       at (5.55,0.92){$\alpha{+}\beta$};
\end{tikzpicture}
\end{center}

\obs{\textbf{Por quê?} O externo é suplementar ao interno adjacente
$\gamma$:\\
externo $= 180° - \gamma = \alpha + \beta$ (pois $\alpha+\beta+\gamma=180°$).}

\end{multicols}

\end{document}
